\section{Результаты}
\begin{table}[t]\caption{Основной сценарий зависимости возвратов: средняя стоимость в тысячах и число выборок, где выбрано сокращение возвратов. Последний столбец --- превышение стоимости планировщика с известным генератором. Тестовые сценарии общие; простое правило и эталон добавлены при последующей диагностике.}\label{tab:primary}\centering\small\setlength{\tabcolsep}{3pt}
\begin{tabular}{p{.42\textwidth}rrr}\toprule
Метод & Стоимость & Возвраты/12 & Превышение\\\midrule
Без действия & 147.8 & 0 & 102.5 \\
Статическое планирование & 92.8 & 0 & 47.5 \\
Отбор по прогнозу & 88.8 & 1 & 43.5 \\
Правило по силе зависимости & 73.0 & 5 & 27.7 \\
Сохранение нагрузочных компонент & 57.4 & 9 & 12.0 \\
Ресурс при 20 & 92.6 & 0 & 47.3 \\
Возвраты при 20 & 45.3 & 12 & 0.0 \\
Известный генератор & 45.3 & 12 & 0.0 \\
\bottomrule\end{tabular}\end{table}
\subsection{Сохранение зависимостей меняет подготовленное действие}
Таблица~\ref{tab:primary} выделяет сценарий зависимости возвратов от нагрузки. Сохранение нагрузочных компонент снижает среднюю стоимость на 38.2\% относительно статического планирования и на 38.1\% относительно добавления ресурса при 20. Соответствующие парные 97.5\% интервалы разности с поправкой на множественные сравнения составляют $[-51.56,-19.40]$ и $[-51.23,-19.32]$ тысячи. Относительно отбора компонент по прогнозу снижение равно 35.4\%; описательный 95\% интервал разности --- $[-46.30,-16.66]$ тысячи. Во всех сравнениях единицей неопределённости служит независимая обучающая выборка.

Анализ действий объясняет различие. Отбор по прогнозу сохраняет зависимость возвратов от нагрузки лишь в 2 из 12 выборок, а соответствующий планировщик выбирает сокращение возвратов в 1. Сохранение нагрузочных компонент выбирает его в 9. Во всех 12 случаях выбран порог 20. Следовательно, наблюдаемый выигрыш относится к выбору действия; преимущества оптимизации времени активации он не доказывает. Каталог действий и поиск по симуляциям одинаковы для всех моделей.

Фиксированное сокращение возвратов при 20 --- необходимое сравнение: средняя стоимость равна 45.3 тысячи против 57.4 при сохранении нагрузочных компонент. Парная разность составляет $12.04$ тысячи с описательным 95\% интервалом $[-1.87,25.94]$. Превосходство над этим правилом не установлено. Правило по силе зависимостей использует ту же обученную нагрузочную модель без сравнения действий в симуляции и получает 73.0 тысячи, выбирая сокращение возвратов в 5 из 12 выборок. Интервал разности с ним, $[-38.90,7.65]$ тысячи, также не позволяет установить превосходство. Эти сравнения связывают результат с выбором модели поведения; необходимости имитационного поиска для данной задачи с двумя действиями они не доказывают.

\subsection{Источник ошибки решения}
Планировщик с известным генератором выбирает сокращение возвратов при 20 во всех выборках основного сценария. Его средняя стоимость на независимом тесте поэтому равна 45.3 тысячи; здесь ретроспективно лучшее фиксированное правило совпадает с этим выбором. Среднее превышение стоимости при сохранении нагрузочных компонент относительно такого эталона знания модели составляет 12.04 тысячи. Всё превышение приходится на 3 выборки с выбором ресурса. В них прогноз стоимости ресурса ниже прогноза для сокращения возвратов, но фактическое ранжирование меняется на обратное (табл.~\ref{tab:choices}). Максимумы нагрузки при калибровке равны лишь 12--19 незавершённым задачам. Это ошибка ранжирования действий при экстраполяции, а не свидетельство другого оптимального порога. Анализ выявляет расхождение решений, не приписывая его однозначно одному коэффициенту и не отождествляя правильный выбор действия с точным восстановлением закона генерации.
\begin{table}[t]\caption{Три выбора ресурса. Разности: ресурс минус сокращение возвратов при 20, в тысячах; отрицательный прогноз означает выбор ресурса.}\label{tab:choices}\centering\small\setlength{\tabcolsep}{3pt}\begin{tabular}{rrrr}\toprule
Seed & Максимум калибровки & Прогноз разности & Фактическая разность\\\midrule
4 & 19 & -1.41 & 51.66 \\
5 & 12 & -0.44 & 42.42 \\
11 & 12 & -4.24 & 50.39 \\
\bottomrule\end{tabular}\end{table}
\subsection{Проверки условий применения}
В исходном тесте возвратов наблюдаемый сигнал предшествует началу тяжёлого события в сопоставленном запуске без действия во всех 289 траекториях риска; среднее опережение равно 108.0. Реакция предотвращает определённое событие в 93 из этих траекторий (32.2\%); снижение стоимости поэтому шире полного предотвращения кризиса. Ни в одной из 192 спокойных тестовых траекторий реакция не активируется. На рис.~\ref{fig:response} показаны траектории для двух отложенных всплесков.

Дополнительная проверка эффекта сохраняет выбранные реакции. При фактическом сокращении возвратов на 30\%, 50\% и 70\% снижение стоимости основного сценария относительно статического планирования составляет 17.0--48.7\%. При фактическом сокращении на 30\% и 70\% описательные парные 95\% интервалы разности стоимости составляют соответственно $[-21.84,-9.67]$ и $[-62.63,-27.85]$ тысячи. Дополнительные материалы содержат каждый механизм, метод и величину эффекта. Среди 9 заданных сочетаний горизонта и стоимости конечного остатка средняя разность относительно статического планирования остаётся благоприятной в 9 случаях. Чувствительность к определению события показана отдельно от стоимости, поэтому выигрыш стоимости не интерпретируется как гарантия доступности.

Выявление структуры только по обучающим данным восстанавливает последовательный видимый язык $A(BC)^+$ во всех 48 журналах. Компилятор исчерпывающе проверяет этот язык по достижимым состояниям автомата и проверяет сохранение одним переходом последовательного исполнения с одним токеном. Новый запуск на заранее выбранной первой выборке возвратов точно воспроизводит выбор правил, стоимость, включение, события и счётчики всех 252 запусков планирования и 96 запусков действий на независимом тесте. Средние времена завершения совпадают с точностью машинного округления. Следовательно, в данном восстановленном последовательном семействе реакция получается через выявленную структуру. Мощности ресурсов и эффекты действий остаются заданными входами; эта проверка не доказывает устойчивости к неточным параллельным структурам.
\begin{figure}[t]\centering\includegraphics[width=\textwidth]{generated/e7_response_ru.pdf}\caption{Реакции в сценарии возвратов при двух отложенных всплесках. Каждая линия усредняет восемь траекторий внутри каждой из 12 выборок, затем выборки. Полоса --- одно стандартное отклонение средних по выборкам. Серые области отмечают всплески; горизонтальная линия --- WIP 100.}\label{fig:response}\end{figure}
